\textbf{Persistent Segment Tree}
\begin{lstlisting}
struct PST {
  struct Node { int l, r, val; };
  vector<Node> t;
  int T = 0;
  PST(int max_nodes) { // estimate 20 * n for
  // max_nodes
    t.resize(max_nodes);
  }
  int build(int b, int e) {
    int cur = T++;
    if (b == e) return cur;
    int m = (b + e) / 2;
    t[cur].l = build(b, m);
    t[cur].r = build(m + 1, e);
    return cur;
  }
  int update(int pre, int b, int e, int i, int v) {
    int cur = T++;
    t[cur] = t[pre];
    if (b == e) {
      t[cur].val += v;
      return cur;
    }
    int m = (b + e) / 2;
    if (i <= m) t[cur].l = update(t[pre].l, b, m, i,
    v);
    else t[cur].r = update(t[pre].r, m + 1, e, i, v);
    t[cur].val = t[t[cur].l].val + t[t[cur].r].val;
    return cur;
  }
  int query(int pre, int cur, int b, int e, int k) {
    if (b == e) return b;
    int cnt = t[t[cur].l].val - t[t[pre].l].val;
    int m = (b + e) / 2;
    if (cnt >= k) return query(t[pre].l, t[cur].l, b,
    m, k);
    return query(t[pre].r, t[cur].r, m + 1, e, k - cnt);
  }
};
int main() {
  ios::sync_with_stdio(0); cin.tie(0);
  int n, q;
  cin >> n >> q;
  vector<int> a(n + 1), root(n + 1), V;
  map<int, int> comp;
  for (int i = 1; i <= n; i++) cin >> a[i],
  comp[a[i]];
  int id = 0;
  V.resize(comp.size() + 1);
  for (auto& [x, _] : comp) comp[x] = ++id, V[id]
  = x;
  PST pst(20 * (n + 1)); // Allocate enough space
  root[0] = pst.build(1, id);
  for (int i = 1; i <= n; i++)
  root[i] = pst.update(root[i - 1], 1, id,
  comp[a[i]], 1);
  while (q--) {
    int l, r, k;
    cin >> l >> r >> k;
    cout << V[pst.query(root[l - 1], root[r], 1, id,
    k)] << '\n';
  }
  return 0;
}
\end{lstlisting}
%the code returns k-th number in a range l to r if
%the range were sorted
\begin{lstlisting}
int V[N], root[N], a[N];
int32_t main() {
  map<int, int>mp;
  int n, q;
  cin >> n >> q;
  for(int i = 1; i <= n; i++) cin >> a[i], mp[a[i]];
  int c = 0;
  for(auto x : mp) mp[x.first] = ++c, V[c] = x.first;
  root[0] = t.build(1, n);
  for(int i = 1; i <= n; i++) {
    root[i] = t.upd(root[i - 1], 1, n, mp[a[i]], 1);
  }
  while(q--) {
    int l, r, k;
    cin >> l >> r >> k;
    cout << V[t.query(root[l - 1], root[r], 1, n, k)]
    << '\n';
  }
  return 0;
}
\end{lstlisting}

\switchcolumn

\textbf{Fenwick Tree}
Point update and range query.
\begin{lstlisting}
template<typename T>
struct BIT {
  int n;
  vector<T> bit;
  BIT(int n = 0) { init(n); }
  void init(int n_) { n = n_; bit.assign(n + 1, 0); }
  void add(int i, T v) { for (; i <= n; i += i & -i)
    bit[i] += v; }
  T sum(int i) { T r = 0; for (; i > 0; i -= i & -i) r
    += bit[i]; return r; }
  T sum(int l, int r) { return sum(r) - sum(l - 1); }
};
\end{lstlisting}
Usage: \texttt{BIT<ll> fenw(5);}

\textit{//Fenwick 2D}
\begin{lstlisting}
struct BIT2D {
  int n, m;
  vector<vector<ll>> bit;
  BIT2D(int n, int m): n(n), m(m) {
    bit.assign(n + 1, vector<ll>(m + 1, 0));
  }
  void add(int x, int y, ll val) {
    for (int i = x; i <= n; i += i & -i)
      for (int j = y; j <= m; j += j & -j)
        bit[i][j] += val;
  }
  ll sum(int x, int y) {
    ll res = 0;
    for (int i = x; i > 0; i -= i & -i)
      for (int j = y; j > 0; j -= j & -j)
        res += bit[i][j];
    return res;
  }
  ll query(int x1, int y1, int x2, int y2) {
    return sum(x2, y2) - sum(x1 - 1, y2) - sum(x2,
    y1 - 1) + sum(x1 - 1, y1 - 1);
  }
};
\end{lstlisting}
\begin{lstlisting}
BIT2D fenw2(4, 4);
fenw2.add(1, 1, 5);
fenw2.add(3, 2, 7);
cout << fenw2.query(1, 1, 3, 3) << "\n"; // sum
// inside rectangle [(1,1)-(3,3)]
\end{lstlisting}

Range update and range query.
\begin{lstlisting}
using ll = long long;
ll bit1[200005] ,bit2[200005],n,m;
void update(ll bit[],ll i,ll val){
  for(;i<=n;i+=(i&(-i))){
    bit[i]+=val;
  }
}
void range_update(ll l, ll r,ll val){
  update(bit1,l,val);
  update(bit1,r+1,-val);
  update(bit2,l,val*(l-1));
  update(bit2,r+1,-(r*val));
}
ll sum(ll bit[],ll i){
  ll re = 0;
  for(;i>0;i-=(i&(-i))){
    re+=bit[i];
  }
  return re;
}
ll prefix_sum(ll i){
  return sum(bit1,i)*i-sum(bit2,i);
}
ll range_sum(ll l,ll r){
  return prefix_sum(r)-prefix_sum(l-1);
}
\end{lstlisting}
\texttt{range\_update(l,r,val);}\\
\texttt{range\_sum(l,r);}

